\chapter*{RESUMEN}
\addcontentsline{toc}{chapter}{RESUMEN}

El proyecto diseñó una máquina para limpiar rodamientos de gran tamaño de rodillos de molienda de alta presión (HPGR) en talleres de mantenimiento minero. Se tomó como referencia un rodamiento de 3032 kg, 1420 mm de diámetro exterior y 515 mm de ancho. Con el modelo de desarrollo de productos VDI 2221 se definieron catorce especificaciones y se generaron cuatro alternativas. Una matriz de decisión seleccionó una configuración en la que el rodamiento permanece inmóvil sobre una mesa de tres apoyos, trasladada por una doble cadena, mientras un distribuidor giratorio en U y una lanza interior aplican el fluido con 24 boquillas. El soporte se verificó con un análisis estructural de barras, un modelo global con torsión no uniforme, elementos finitos de los centradores y el diseño de sus uniones soldadas. La red hidráulica se resolvió por aproximaciones sucesivas, y el mojado se simuló con un modelo de fase discreta verificado con tres mallas. El bastidor alcanzó una razón demanda/capacidad máxima de 0,82 y el error de colocación del rodamiento quedó en 4,24 mm frente a un límite de 5 mm. La red demandó 221 L/min a unos 53 m, con presiones de boquilla de 4,00 a 4,13 bar. Un par de boquillas de abanico estrecho multiplicó por 4,8 la dosis en la zona exterior de la cara inferior, la superficie crítica, y redujo de 37,6 a 14,1 min el tiempo de lavado estimado de esa cara, que es la que fija el tiempo de la máquina. El presupuesto de inversión se estimó en S/~314\,310 sin IGV. El diseño es viable como ingeniería preliminar verificada por simulación y cumple trece de las catorce especificaciones en los modelos; su eficacia de limpieza y su justificación económica requieren ensayos y datos que están fuera del alcance del proyecto.

\medskip
\noindent\textbf{Palabras clave:} \textit{rodamiento HPGR; limpieza por aspersión; diseño de máquinas; dinámica de fluidos computacional; elementos finitos; mantenimiento minero.}
